% ToolFrame 6.4 RA-L / ICRA paper outline draft.
% RA-L official guidance points authors to the IEEE journal LaTeX template
% (IEEEtran). If submitting as an ICRA conference paper instead, port the body
% to the official ICRA ieeeconf/IEEE conference template and keep the same
% section budget.
\documentclass[journal]{IEEEtran}

\usepackage{amsmath,amssymb}
\usepackage{graphicx}
\usepackage{booktabs}
\usepackage{multirow}
\usepackage{xcolor}

\newcommand{\todo}[1]{\textcolor{red}{TODO: #1}}
\newcommand{\method}{ToolFrame}

\title{ToolFrame: Let the Tool Define the Frame for Dexterous Manipulation}

\author{Author~Names%
\thanks{Manuscript received: \todo{date}; Revised: \todo{date}; Accepted: \todo{date}.}%
\thanks{This paper was recommended for publication by \todo{Editor name} upon evaluation of the Associate Editor and Reviewers' comments.}%
\thanks{Authors are with \todo{affiliation}. Correspondence: \todo{email}.}%
}

\begin{document}
\maketitle

\section*{Paper Positioning Notes}
\textbf{Candidate titles.}
\begin{itemize}
  \item \method: Let the Tool Define the Frame for Dexterous Manipulation
  \item \method: Tool-First Visual Grounding for Dexterous Manipulation
  \item \method: Tool-Centric Visual Grounding for Dexterous Tool Use
  \item \method: Compiling Tool-Centric Visual Evidence into Dexterous Manipulation Subgoals
  \item Planning Dexterous Tool Use Around Functional Tool Geometry
\end{itemize}
\textbf{One-sentence pitch.} \method{} makes the tool, rather than the hand, the
organizing frame for dexterous manipulation: the VLM selects public tool and
scene evidence, while a geometry compiler produces whole-tool $SE(3)$ subgoals
that the hand follows.

\begin{abstract}
Dexterous tool use is defined by the tool: which part should contact the scene,
which axis should align with the intended action, and where the resulting
whole-tool motion should go. Recent
vision-language models (VLMs) offer strong semantic grounding, but directly
using them to generate robot trajectories can entangle image interpretation,
geometry, controller feasibility, and hidden benchmark assumptions. We propose
\method{}, an image-grounded tool-use pipeline that lets the tool define the
planning frame while keeping visual grounding and robot execution separate:
a VLM selects sparse public scene references and public tool functional
references, while a geometric compiler converts the references into whole-tool
$SE(3)$ subgoals for fixed-goal policy execution. The method consumes natural
task language, public RGB-D scene evidence, public tool meshes, and
reviewed-or-pending tool-part evidence, with final headline claims restricted
to contracts whose review status is explicitly marked. It excludes hidden
evaluator references, oracle
target axes, oracle trajectories, task-family fallback, primitive fallback,
repair-loop fallback, and seed-specific shortcuts. \method{} grounds
start/end/path scene key-points through RGB-D, binds tool.contact\_point and
tool.motion\_axis to those references, compiles feasible whole-tool subgoals,
and executes them with a SimToolReal fixed-goal policy. We evaluate with task
contracts that report final task success and attribute failures to functional
grounding, scene grounding, subgoal validity, policy execution, or task-effect
failure. Final headline results should be filled only from frozen
launcher\_summary files: \todo{insert final N, task success rates, and failure
attribution after release rerun}. Candidate and diagnostic tasks stress
insertion, striking, and side-board marker tracing without promoting unreviewed
contracts to headline statistics.
\end{abstract}

\begin{IEEEkeywords}
Robot learning, tool use, visual grounding, vision-language models, motion
planning, benchmark evaluation.
\end{IEEEkeywords}

\section{Introduction}
\textbf{Opening problem.} Dexterous tool use is often planned around the robot
hand, but the task is defined by the tool. A brush has bristles that push, an
eraser has a flat contact patch that wipes, a screwdriver has a tip and an
insertion axis, and a marker has a tip that leaves a trace. A robot must
therefore answer two coupled visual questions before it can act: which part of
the tool is functional, and where in the scene should that function be applied?

\textbf{Why this is not just VLM planning.} VLMs can often point to a target
object or describe an intended action, but a robot still needs metric contact
geometry, tool-frame orientation, and controller-compatible subgoals. If the VLM
is asked to output a trajectory directly, the resulting plan can hide geometry
errors, silently inherit task-family defaults, or depend on evaluator-only
target axes. This is especially problematic for benchmarks, where the line
between method inputs and private success predicates must stay visible.

\textbf{Core idea.} \method{} lets the tool define the planning frame. The VLM
is treated as a sparse evidence selector rather than a trajectory generator.
The scene side is grounded from public RGB-D as start\_point, optional
path\_points, end\_point, and scene.direction. The tool side is grounded from
public part evidence as tool.contact\_point and tool.motion\_axis. A compiler
then solves whole-tool $SE(3)$ subgoals that bind the selected tool references
to the selected scene references. The hand follows these tool-centric fixed
goals, and task contracts evaluate the executed effect.

\textbf{Paper narrative.} Figure~\ref{fig:pipeline} should be the visual anchor:
A) tool-centric prior extraction, B) scene key-point grounding, C)
geometry-constrained optimization, and D) policy rollout and evaluation. The
paper should repeatedly emphasize that the VLM never sees hidden evaluator refs
and never emits robot trajectories.

\textbf{Contributions.} This paper makes three contributions:
\begin{enumerate}
  \item We introduce an auditable VLM interface for robotic tool use that
  restricts model outputs to sparse scene references and functional tool
  references, rather than robot trajectories, oracle axes, or task-family
  primitives.
  \item We present a geometry-constrained compiler that binds lifted RGB-D
  scene references with tool contact-point and motion-axis references to
  produce executable whole-tool $SE(3)$ subgoals.
  \item We define a task-contract evaluation protocol that reports final task
  success together with module-level failure attribution, while separating
  reviewed-or-explicitly-marked headline tasks from provisional candidate and
  diagnostic tasks.
\end{enumerate}

\section{Related Work}
\subsection{Language and Vision-Language Models for Robot Planning}
Write this subsection as a contrast between semantic planning and metric
execution. Cover VLM/LLM robot planners, code-as-policies style systems,
affordance prompting, and image-conditioned action proposal. The key distinction
is that \method{} uses the VLM only for auditable sparse references, not for
robot trajectories, hidden target axes, or task-family primitives.

\subsection{Tool-Use Affordances and Functional Part Grounding}
Cover tool affordance learning, part segmentation, functional region
prediction, contact-point prediction, and object-part representations. Position
\method{} as tool-centric: it converts segmented part evidence into a small
functional reference interface, tool.contact\_point and tool.motion\_axis, that
can be inspected and reused by the compiler.

\subsection{Geometric Motion Planning and Constraint-Based Manipulation}
Discuss task-space constraints, contact-aware manipulation, keyframe/subgoal
planning, and $SE(3)$ pose optimization. The packaging point is that \method{}
does not introduce named manipulation primitives such as sweep, strike, or
scrape. Those are task-language interpretations; the compiler sees a compact
constraint vocabulary: point\_match, point\_offset, axis\_align, move\_along,
orientation\_offset\_about\_axis, and clearance.

\subsection{Robotic Manipulation Benchmarks and Evaluation Protocols}
Compare with benchmarks that report final task success only. Motivate
contract-based success and module-level attribution while preserving hidden
evaluator data as evaluator-only information.

\section{Method}
\begin{figure*}[t]
\centering
\fbox{\parbox{0.92\textwidth}{\centering
Placeholder for four-panel method figure. A: tool-part evidence to
tool.contact\_point/tool.motion\_axis. B: RGB-D scene image to sparse
scene refs. C: geometry-constrained whole-tool $SE(3)$ optimizer. D:
SimToolReal fixed-goal rollout and task-contract evaluation.}}
\caption{\method{} separates semantic grounding from robot trajectory
generation. VLMs select sparse public evidence; the compiler produces
whole-tool subgoals; the policy executes fixed goals; task contracts score
executed effects.}
\label{fig:pipeline}
\end{figure*}

\subsection{Problem Setup and Method Boundary}
Inputs: task language, calibrated public RGB-D scene evidence, public tool mesh
and tool-part evidence, and optional robot/gripper state. Private simulator
object states, hidden evaluator masks, oracle contact annotations, and success
predicates are evaluation-only.

Forbidden method inputs: hidden evaluator refs, oracle target axes, oracle
trajectories, tool-family fallback, task-family fallback, primitive fallback,
repair-loop fallback, automatic global-axis retries, and seed-specific
shortcuts.

\textbf{Representation.} The method uses two public reference sets:
\begin{align}
\mathcal{R}_{tool} &= \{p^{contact}_{tool}, a^{motion}_{tool}\}, \\
\mathcal{R}_{scene} &= \{p^{start}_{scene}, p^{path}_{scene}, p^{end}_{scene},
a^{direction}_{scene}\}.
\end{align}
The compiler is responsible for converting these references into executable
whole-tool poses. This representation should be introduced early because it is
the paper's organizing abstraction.

\subsection{A. Tool-Centric Prior Extraction}
\textbf{Input.} Tool mesh, P3SAM/SAMPart3D part evidence, task language, and
reviewed visual evidence cards.

\textbf{Process.} The VLM first selects a functional part from segmented part
evidence. Only the selected part is then shown for contact-point and motion-axis
selection, with part bounding boxes and face highlights available for review.

\textbf{Output.} Public tool refs under configs/tool\_refs/\textless tool\_id\textgreater:
part\_selection, contact\_point\_selection, axis\_binding, and tool\_part\_graph.
These produce only tool.contact\_point and tool.motion\_axis; the method does
not invent coordinates or use evaluator labels as method inputs.

\textbf{What to show.} Include a compact evidence figure with the selected
functional part, the selected contact face/point, and the selected motion-axis
face or marker. The caption should state that evaluator-side tool labels are
used only for audit/failure attribution, not for method selection.

\subsection{B. Scene Key-Points Grounding}
\textbf{Input.} Public RGB-D images, depth, camera calibration, and natural task
language.

\textbf{VLM role.} Mark sparse scene refs: start\_point, optional ordered
path\_points, end\_point, and scene.direction when needed. For tasks without a
visible endpoint, the scene direction still comes from public VLM scene
grounding or an explicit task contract; no hidden target center replaces the
VLM pixel.

\textbf{Lifting.} Pixels are lifted through RGB-D and camera calibration to
3D scene refs. If an endpoint is visible, action distance comes from the lifted
start/end geometry. If not image-separable, the endpoint defaults to the start
point and the action distance is zero unless a public geometric projection is
specified by the contract.

\textbf{What to show.} Include one top-down or side-board grounding overlay per
task family, with visible start/path/end refs. The figure should make clear that
scene refs come from public RGB-D, not simulator target centers.

\subsection{C. Geometry-Constrained Optimizer}
\textbf{IR.} Build staged constraints for grasp/lift, align, approach, and
action. Each method stage contains terminal subgoal constraints and path
constraints; approach is an offset behind the first contact point along the
negative VLM motion direction.

\textbf{Optimization objective.} Solve whole-tool poses $T^i_{world,tool}$ such
that
\begin{align}
T_{world,tool} \, p_{tool}^{contact} &\approx p_{scene}^{ref}, \\
R_{world,tool} \, a_{tool}^{motion} &\parallel a_{scene}^{direction}.
\end{align}
The compiler lowers the constraint program to ToolPoseSubgoals and then to
simtoolreal\_fixed\_goals.json.

\textbf{Boundary.} Missing staged axis constraints are errors. The compiler does
not reselect tool faces, guess brush/hammer defaults, snap to hidden simulator
targets, or repair failed seeds with private references.

\textbf{Compiler outputs.} Report both the symbolic constraint program and the
lowered SimToolReal fixed goals as artifacts. The paper should describe this as
a compiler boundary rather than a learned policy boundary: VLM refs are not
backend commands, and backend fixed-goal fields are not task primitives.

\subsection{D. Policy Rollout and Evaluation}
SimToolReal executes the fixed goals. The policy receives compiled subgoals, not
task-family primitives or VLM trajectories. The evaluator applies task
contracts after rollout and writes rollout summaries, visual evidence manifests,
task-success checks, and failure attribution.

\textbf{Evaluation interface.} Task success predicates are evaluator-side
contracts. They may inspect rollout effects, contact coverage, insertion depth,
or produced ink evidence, but they never feed back into the method path.

\section{Task-Contract Benchmark}
\subsection{Task Tiers}
\begin{table*}[t]
\centering
\caption{ToolFrame 6.4 task tiers. Only current/e2e tasks should be headline
statistics unless a review note closes the provisional contract.}
\begin{tabular}{p{0.13\textwidth}p{0.60\textwidth}p{0.20\textwidth}}
\toprule
Tier & Tasks & Reporting role \\
\midrule
Headline & brush\_sweep/forward; eraser\_wipe/tabletop\_wipe & final table after frozen release \\
Candidate & brush\_sweep/left; hammer\_strike/down; screwdriver\_insert/wide\_socket; marker\_trace/write\_l,c & provisional evidence \\
Diagnostic & screwdriver\_insert/deep\_socket; dense variants & stress only \\
Draft & side-board eraser/hammer/screwdriver & smoke and contract design \\
\bottomrule
\end{tabular}
\end{table*}

\subsection{Task Contracts}
Brush sweep succeeds by final target-block occupancy inside the visible red
outline zone; directional displacement is diagnostic. Eraser wipe succeeds by a
diagnostic path gate, lifted refs on the visible raised tabletop HDU board, and
flat-eraser contact coverage at least 30\% of the target long-axis extent.
Hammer strike is provisional and uses VLM strike direction plus semantic
contact/strike-direction gates.
Screwdriver wide-socket insertion is provisional and uses tip radial alignment
and measured insertion depth. Marker tracing is a hard side-board candidate
judged by produced marker-tip ink evidence; contact/tracking are diagnostics.

\subsection{Metrics and Failure Attribution}
The main reporting table uses one row per headline task:
\begin{quote}
task, contract status, review note, $N$, functional grounding rate, scene
grounding rate, subgoal validity rate, policy execution rate, task success rate.
\end{quote}
Task success is the final task metric. The intermediate rates identify the
first stage where public evidence or execution artifacts fail. Missing
artifacts count as failures at the first missing evidence stage; ambiguous
evidence is reported as needs\_review.

\begin{table*}[t]
\centering
\caption{Headline benchmark results. Rows are included only when the task
contract is current/e2e or its pending labels are explicitly marked. Fill
only from frozen launcher\_summary files and matching per-task summary.json
files under reports/paper\_release\_6.4/.}
\scriptsize
\setlength{\tabcolsep}{3pt}
\begin{tabular}{p{0.17\textwidth}p{0.10\textwidth}p{0.22\textwidth}cccccc}
\toprule
Task & Status & Review note & N & Func. & Scene & Subgoal & Policy & Success \\
\midrule
brush\_sweep/forward & provisional & success predicate reviewed; red-brush labels pending & \todo{} & \todo{} & \todo{} & \todo{} & \todo{} & \todo{} \\
eraser\_wipe/tabletop\_wipe & reviewed & current contract; flat-eraser labels approved & \todo{} & \todo{} & \todo{} & \todo{} & \todo{} & \todo{} \\
\bottomrule
\end{tabular}
\end{table*}

\begin{table*}[t]
\centering
\caption{Candidate and diagnostic evidence. These rows stress additional task
families and report readiness/failure attribution only; they are not included
in headline aggregate statistics.}
\begin{tabular}{p{0.23\textwidth}p{0.13\textwidth}p{0.53\textwidth}}
\toprule
Task & Status & Reporting note \\
\midrule
brush\_sweep/left & provisional & Same directional-displacement family as forward; needs review before headline statistics. \\
hammer\_strike/down & provisional & Compile plus semantic contact/strike-direction gate; hammer labels pending. \\
screwdriver\_insert/wide\_socket & provisional & Tip radial alignment and insertion-depth checks; long-screwdriver labels pending. \\
marker\_trace/write\_l,c & provisional & Hard side-board candidates judged by produced marker-tip ink; marker labels pending. \\
screwdriver\_insert/deep\_socket and dense variants & stress & Diagnostic insertion stress splits; excluded from default headline statistics. \\
side-board eraser/hammer/screwdriver & draft & Smoke and contract-design tasks only. \\
\bottomrule
\end{tabular}
\end{table*}

\textbf{Artifact boundary.} Final headline statistics are populated only from
completed frozen launcher\_summary files under reports/paper\_release\_6.4/ and
the corresponding per-task summary.json files. Smoke runs and failed
render/debug launchers are retained as provenance, not as benchmark results.
reports/mainline\_6.4/ is a historical audit/migration snapshot and must not be
reported as final \method{} 6.4 success statistics.

\section{Experiments}
\subsection{Experimental Questions}
Q1: Does public image-grounded evidence suffice to produce valid whole-tool
subgoals without hidden evaluator refs? Q2: Which stage limits end-to-end task
success? Q3: How do reviewed headline tasks differ from provisional candidate
and diagnostic tasks? Q4: What failures become attributable when VLM grounding
is separated from geometric compilation, policy execution, and task-effect
evaluation?

\subsection{Comparisons}
The main controlled comparisons should be boundary-preserving ablations. We do
not treat legacy family defaults, primitive fallbacks, or repair loops as
\method{} baselines, because they violate the evidence boundary that the paper
is designed to study. If reported at all, historical family-default behavior
belongs in an appendix provenance analysis, never as clean-method evidence.

\begin{itemize}
  \item Public human-clicked scene refs: replace VLM scene pixels with
  human-clicked pixels on the same public RGB-D views. Report only as an
  upper-bound scene-grounding diagnostic; no simulator target centers, hidden
  masks, or evaluator refs are used.
  \item Tool-reference diagnostic: hold scene refs fixed and vary public
  tool.contact\_point/tool.motion\_axis choices to measure functional grounding
  sensitivity. \todo{define reviewed choices before reporting}
  \item VLM scene-ref diagnostic: compare successful and failed VLM sparse refs
  using the same compiler/evaluator path to isolate scene marking and lifting
  errors. \todo{fill from per-trial artifacts}
\end{itemize}

\subsection{Ablation Design}
\begin{table}[t]
\centering
\caption{Suggested ablations and diagnostics. All preserve the no-hidden-ref
boundary and should be reported separately from final task success.}
\begin{tabular}{p{0.33\columnwidth}p{0.57\columnwidth}}
\toprule
Ablation & Expected diagnostic \\
\midrule
Remove tool motion axis & compile/preflight failure; verifies that the compiler does not silently guess task-family axes \\
Remove selected contact point & compile or contact/action failure; verifies that whole-tool goals require an explicit functional point \\
2D refs without RGB-D lifting & scene/subgoal failures \\
No path points for marker C & semantic ink recognition drop \\
Strict tracking diagnostic & reports disagreement between pose tracking and task-effect success; does not replace the task contract \\
Candidate vs headline contracts & readiness and attribution analysis \\
\bottomrule
\end{tabular}
\end{table}

\subsection{Figures}
Fig. 1: Four-panel pipeline (A tool prior, B scene key-points, C optimizer,
D rollout/evaluation). Fig. 2: tool evidence cards and selected
contact/axis refs. Fig. 3: scene RGB-D grounding overlays. Fig. 4: compiled
whole-tool subgoals over rollout frames. Fig. 5: failure attribution stacked
bar by task.

\section{Limitations}
\method{} depends on visible RGB-D evidence and calibrated lifting; occlusion,
poor depth, or ambiguous images can block scene grounding. Tool refs require
public part evidence and, for final labels, human review; this limits immediate
coverage of new tools. The fixed-goal SimToolReal policy can fail even when
compiled subgoals are valid. Current candidate tasks such as screwdriver
insertion and marker tracing need additional contract review before headline
statistics. The method does not yet learn from failures online and
intentionally excludes repair-loop fallback.

\section{Conclusion}
\method{} shows that robotic tool use can be structured as evidence-grounded
compilation rather than direct trajectory generation. By restricting VLMs to
sparse, auditable scene and tool references, and by compiling those references
into whole-tool $SE(3)$ subgoals, the system exposes geometric assumptions that
often remain hidden in tool-use pipelines. The task-contract benchmark further
separates final task success from functional grounding, scene grounding,
subgoal validity, policy execution, and task-effect failures. This separation
does not remove the difficulty of contact-rich tool use, but it makes remaining
failures measurable and reviewable. \todo{insert final quantitative takeaway
after release rerun.}

\appendices
\section{Appendix Checklist}
\begin{itemize}
  \item Official RA-L/IEEE template files and compile log.
  \item Full task contract table with thresholds and review status.
  \item Tool-ref evidence manifest for each tool in configs/tool\_refs/.
  \item Scene evidence examples and VLM grounding prompts.
  \item Constraint-program schemas and examples.
  \item Per-task launcher\_summary.{json,csv,md} and per-task summary.json.
  \item Failure attribution definitions and ambiguous needs\_review examples.
  \item Boundary-preserving comparison and ablation configs.
  \item Asset/license statement for SimToolReal and benchmark assets.
  \item Negative-results note for failed or provisional release runs.
\end{itemize}

\section{One-Page Compressed Outline}
\textbf{Title.} ToolFrame: Let the Tool Define the Frame for Dexterous
Manipulation.
\textbf{Claim.} ToolFrame makes the tool, rather than the hand, the organizing
frame for dexterous manipulation: sparse public tool and scene references are
compiled into whole-tool $SE(3)$ goals that the hand follows.
\textbf{Method.} A) extract tool.contact\_point and tool.motion\_axis from
public tool-part evidence; B) ground start/path/end refs and scene.direction
from public RGB-D; C) solve geometry-constrained whole-tool subgoals; D) execute
with SimToolReal fixed goals and score with task contracts.
\textbf{Benchmark.} Headline: brush\_sweep/forward and
eraser\_wipe/tabletop\_wipe after frozen rerun. Candidate: left sweep, hammer,
wide-socket screwdriver, marker L/C. Diagnostic: deep/dense screwdriver. Draft:
side-board eraser/hammer/screwdriver.
\textbf{Main table.} task, contract status, review note, N, functional
grounding, scene grounding, subgoal validity, policy execution, task success.
Fill TODO only from frozen launcher summaries.
\textbf{Comparisons.} public human-clicked scene refs, tool-reference
sensitivity, VLM scene-ref diagnostics, remove tool axis/contact point,
RGB-D-lifting ablation, tracking-vs-effect disagreement.
\textbf{Boundary.} No hidden evaluator refs, oracle axes, oracle trajectories,
family fallbacks, primitive fallbacks, repair-loop fallbacks, or seed shortcuts.

\section{Eight-Page RA-L/ICRA Page Budget}
\begin{table*}[t]
\centering
\caption{Suggested 8-page body allocation.}
\begin{tabular}{p{0.12\textwidth}p{0.22\textwidth}p{0.58\textwidth}}
\toprule
Pages & Section & Content \\
\midrule
0.0--0.5 & Abstract/Intro lead & problem, thesis, contributions \\
0.5--1.4 & Introduction & motivation, boundary, summary figure \\
1.4--2.1 & Related Work & VLM planning, affordances, constraints, benchmarks \\
2.1--2.8 & Setup/Boundary & inputs, forbidden signals, task refs \\
2.8--3.5 & A/B Method & tool prior extraction, scene key-point grounding \\
3.5--4.5 & C/D Method & optimizer, fixed-goal rollout, evaluator contracts \\
4.5--5.2 & Task-Contract Benchmark & task tiers, contracts, metrics \\
5.2--6.4 & Experiments & main table, module attribution, qualitative rollouts \\
6.4--7.2 & Comparisons/Ablations & boundary-preserving diagnostics and TODO results \\
7.2--7.7 & Limitations & visibility, tool-ref review, policy limits \\
7.7--8.0 & Conclusion & final takeaway and release statement \\
\bottomrule
\end{tabular}
\end{table*}

\bibliographystyle{IEEEtran}
\bibliography{references}

\end{document}
